\documentclass[12pt]{article}
\usepackage[utf8]{inputenc}
\usepackage[T1]{fontenc}
\usepackage{hyperref}
\title{Bu Language Resources for K-12 Curriculum Development}
\author{Compiled by Claude Sonnet 5.5 (claude-sonnet-5-5) for Andr\'as Kornai}
\date{}
\begin{document}
\maketitle

\section{Introduction}
Bu (ISO 639-3 zbu; also Buu, Zaranda, Zarandawa) is a Chadic (Afro-Asiatic) language of southern Bauchi State, Nigeria, in the South-Bauchi West group. It is not mutually intelligible with Bolu, Pelu or Geji. The Endangered Languages Project lists it as severely endangered, and it reports that most speakers have shifted to Hausa and Fulfulde. This differs from the reference description in the prompt, which calls the language stable. Only a few speakers are reported to remain in Zaranda Habe. The searches found little material, and nothing for schools.

\section{Primary Educational Resources}
None found. The language is not known to be taught in schools.

\section{Linguistic and Reference Materials}
\begin{itemize}
\item Bernard Caron, ``Bu: aka Zaranda (Chadic, South-Bauchi West): Wordlist and grammatical notes'' (HAL, LLACAN): \url{https://hal-cnrs.archives-ouvertes.fr/LLACAN/halshs-00644360v1}
\item Endangered Languages Project entry for Buu (Nigeria): \url{https://www.endangeredlanguages.com/lang/10444}
\item Endangered Languages Project bibliography for Buu: \url{https://endangeredlanguages.com/lang/10444/bibliography}
\end{itemize}

\section{Religious and Community Resources}
No religious texts were confirmed. Global Recordings Network lists recordings under the names Bu and Zarandawa (see below).

\section{Digital and Audio Resources}
\begin{itemize}
\item Global Recordings Network, Bu (Bauchi State): \url{https://www.globalrecordings.net/en/language/zbu}
\item Global Recordings Network, Bu: Zarandawa: \url{https://globalrecordings.net/en/language/4989}
\end{itemize}

\section{Curriculum Development Context}
Hausa and Fulfulde appear to be replacing Bu in daily life, so any curriculum use of Bu would be revitalization work. It would rest on Caron's wordlist and grammatical notes and on new recordings made with the remaining speakers. Nothing in the search results shows a standard orthography or school materials.

\end{document}